\section{Measurement framework}
\label{sec:methods}

The paper's quantitative content is an accounting decomposition rather than an
estimator, and we state it precisely so that the numbers in
Section~\ref{sec:results} can be reproduced or disputed.

\subsection{The benchmark and the residual}

For year $t$, let $g^{MW}_t$ be the growth of the statutory minimum wage,
$g^{A}_t$ the growth of real output per person employed, and $\pi_t$ consumer
price inflation. Define the benchmark nominal wage growth and the policy
residual as
%
\begin{equation}
  g^{\text{benchmark}}_{t} = (1 + g^{A}_{t})(1 + \pi_{t-1}) - 1,
  \qquad
  r_{t} = g^{MW}_{t} - g^{\text{benchmark}}_{t}.
  \label{eq:residual}
\end{equation}
%
The benchmark answers a specific question: by how much would the wage floor have
had to rise to keep pace with productivity while compensating for the previous
year's inflation? The residual is the part of statutory growth that did neither.

Two conventions matter. The benchmark \emph{compounds} rather than adds. At the
inflation rates Portugal experienced in the early 1980s the additive
approximation understates it by more than a percentage point, which exceeds most
of the residuals being measured. And inflation enters lagged, because a wage set
in December for the following January can respond to inflation already observed
but not to inflation not yet realised.

\subsection{Annual conventions}

Acts take effect on a stated date, which is 1 January in most years but not all:
the wage changed mid-year in 1975, 1978, 1979, 1980, 1981, 1989 and 2014. Two
annual series are therefore constructed, one giving the level in force on
1~January and one averaging the level over the days of the year. The results
below use the day-weighted average, which is the convention that matches annual
national-accounts aggregates. The two agree in every year without a mid-year
act.

\subsection{Dynamic specification}

Where dynamic responses are estimated in Section~\ref{sec:robustness}, they are
traced by local projections \citep{jorda2005projections}, which estimate the
response at each horizon by a separate regression rather than extrapolating it
from a fitted process.

\subsection{What the residual is not}

The residual is descriptive. A positive value means the wage floor rose faster
than the benchmark; it does not mean the increase caused inflation, and the
correlation between the residual and subsequent inflation cannot settle that
question, because statutory changes respond to the same conditions that drive
prices. Section~\ref{sec:robustness} sets out what would be needed to make a
causal statement, and why the Portuguese case does not currently supply it.
